% !TEX program = xelatex
% !BIB TS-program = biber

% **************************************************
% Document Class Definition
% **************************************************
\documentclass[%
    paper=A4,               % paper size --> A4 is default in Germany
    twoside=true,           % onesite or twoside printing
    openany,              % doublepage cleaning ends up right side
    parskip=full,           % spacing value / method for paragraphs
    chapterprefix=true,     % prefix for chapter marks
    11pt,                   % font size
    headings=normal,        % size of headings
    bibliography=totoc,     % include bib in toc
    listof=totoc,           % include listof entries in toc
    titlepage=on,           % own page for each title page
    captions=tableabove,    % display table captions above the float env
    draft=false,            % value for draft version
]{scrreprt}

\newcommand{\reportTitle}{\textcolor{ugent_blue}{Your report title}} 
\newcommand{\reportName}{\textbf{Your name}}
\newcommand{\reportSubject}{\emph{Your report subject}}
\newcommand{\reportDate}{}
\newcommand{\reportVersion}{}

\usepackage[english]{babel} % babel system, adjust the language of the content
\PassOptionsToPackage{% setup clean thesis style
    figuresep=colon,%
    sansserif=true,%
    hangfigurecaption=false,%
    hangsection=true,%
    hangsubsection=true,%
    colorize=full,%
    colortheme=bluemagenta,%
    bibsys=biber,%
    bibfile=notes,%
    bibstyle=numeric-comp,%
    wrapfooter=false,%
}{cleanthesis}
\usepackage{cleanthesis}

\hypersetup{% setup the hyperref-package options
    pdftitle={\reportTitle},    %   - title (PDF meta)
    pdfsubject={\reportSubject},%   - subject (PDF meta)
    pdfauthor={\reportName},    %   - author (PDF meta)
    plainpages=false,           %   -
    colorlinks=false,           %   - colorize links?
    pdfborder={0 0 0},          %   -
    breaklinks=true,            %   - allow line break inside links
    bookmarksnumbered=true,     %
    bookmarksopen=true          %
}

\definecolor{ugent_blue}{RGB}{30, 100, 200}
\definecolor{sciences}{RGB}{45, 140, 168}

\usepackage{algorithm}              % http://ctan.org/pkg/algorithms
\usepackage{algpseudocode}          % http://ctan.org/pkg/algorithmicx
\usepackage{amsmath}
\usepackage{amsfonts}             % mathfrak
\usepackage{amssymb}              % \mathbb
\usepackage[mathscr]{euscript}    % \mathscr
\usepackage{chemmacros}
\usepackage{float}
\usepackage{graphicx}
\usepackage{interval}
\usepackage{mathtools}            % dcases, DeclareMathOperator, DeclairePairedDelimiter
\usepackage{physics}               % physics typesetting, especially bra-ket notation
\usepackage{siunitx}
\usepackage{subcaption}
\usepackage{tensor}
\usepackage{xcolor}
\usepackage{ifxetex}
\usepackage{multicol}
\usepackage{pdfpages}
\usepackage{epigraph}
\usepackage{bm}
\usepackage[normalem]{ulem}
\usepackage[version=3]{mhchem}
\usepackage[linewidth=1pt]{mdframed}
\usepackage{enumitem}


\DeclareSymbolFont{euleroperators}{U}{eur}{m}{n}
\SetSymbolFont{euleroperators}{bold}{U}{eur}{b}{n}

\makeatletter
\renewcommand{\operator@font}{\mathgroup\symeuleroperators}
\makeatother

% \usepackage[utf8]{inputenc} 


\ifxetex
  \usepackage{fontspec}
	% \setmainfont[
	% Path = ./,
	% BoldFont = UGentPannoText-SemiBold.otf,
	% ItalicFont = UGentPannoText-SemiLight.otf,
	% BoldItalicFont  = UGentPannoText-Medium.otf]
	% {UGentPannoText-Normal.otf}
	\newfontfamily\PannoSemiLight{UGentPannoText-SemiLight.otf}
	\newfontfamily\PannoMedium{UGentPannoText-Medium.otf}
	\newfontfamily\PannoSemiBold{UGentPannoText-SemiBold.otf}
	\newcommand{\PannoLarge}{\LARGE}
    \newcommand{\PannoNormalsize}{\large}
    \renewcommand{\helv}{\fontspec{UGentPannoText-Normal.otf}\fontsize{9}{11}\selectfont}
    \renewcommand{\book}{\fontspec{UGentPannoText-Normal.otf}\fontsize{11}{13}\selectfont}
    \renewcommand{\tgherosfont}{\fontspec{UGentPannoText-Normal.otf}}
    \renewcommand{\thesischapterfont}{\color{ctcolorblack}\huge\fontspec{UGentPannoText-Normal.otf}}
\else
  % \usepackage[utf8]{inputenc}           % Om niet ascii karakters rechtstreeks te kunnen typen
	\usepackage{lmodern}
	\usepackage[T1]{fontenc}
	\renewcommand{\familydefault}{\sfdefault}
	\newcommand{\PannoSemiLight}{}
	\newcommand{\PannoMedium}{}
	\newcommand{\PannoSemiBold}{}
	\newcommand{\PannoLarge}{\Large}
	\newcommand{\PannoNormalsize}{\normalsize}
\fi

\graphicspath{{img/}}
\linespread{1.2}
\numberwithin{equation}{section}
\renewcommand{\thesection}{\arabic{section}}

\title{\reportTitle}
\author{\reportName\\GQCG}

% **************************************************
% Document CONTENT
% **************************************************
\begin{document}

    \begin{titlepage}
        \hspace{-1cm}\includegraphics[height=3cm]{faculty.png}
        \begin{center}
            \vspace{2cm}
            \Huge{\reportTitle} \\
            \vspace{1cm}
            \Large{\reportName} \\
            \vspace{2cm}
            \includegraphics[height=2cm]{GQCG.png} \\
            \vspace{0.5cm}
            \large{Ghent University} \\ 
            \large{\today}
        \end{center}
        \vspace{4cm}
        \hspace{-1cm}\includegraphics[height=4cm]{UGent.png}
    \end{titlepage}
    % --------------------------
    % rename document parts
    % --------------------------
    \renewcaptionname{english}{\figurename}{Fig.}
    \renewcaptionname{english}{\tablename}{Tab.}

    \setcounter{tocdepth}{2}		% define depth of toc

    % --------------------------
    % Body matter / main matter
    % --------------------------
    \pagenumbering{arabic}			% arabic page numbering
    \setcounter{page}{1}			% set page counter
    \pagestyle{maincontentstyle} 	% fancy header and footer

    % The following input statements correspond to the different chapters of these notes
    \section{Template section}

This is a report template. Citations can be used as via the included \emph{.bib} file \cite{acke2020}.
    \clearpage

    % --------------------------
    % Back matter
    % --------------------------

    {%
    \setstretch{1.1}
    \renewcommand{\bibfont}{\normalfont\small}
    \setlength{\biblabelsep}{0pt}
    \setlength{\bibitemsep}{0.5\baselineskip plus 0.5\baselineskip}
    \printbibliography[nottype=online]
    \printbibliography[heading=subbibliography,title={Webpages},type=online,prefixnumbers={@}]
    }
    
    % **************************************************
    % End of Document CONTENT
    % **************************************************
    \end{document}
